\chapter{Clase Taller 2}

\section{Preguntas Think--Pair--Share Transformaciones de Lorentz}

\textbf{1. Avance de un frente de luz en S}
Un destello se produce en el origen común de los sistemas inerciales $S$ y $S'$ cuando ambos coinciden ($t= t'= 0$). A un tiempo $t$ la onda luminosa se ha expandido como una esfera. ¿Cuál es la coordenada que describe la ubicación del frente de onda del destello según $S$?
\begin{enumerate}[a)]
    \item $x= c/t$
    \item \textcolor{red}{$x= ct$}
    \item $x= \gamma ct$
    \item $x= vt$
\end{enumerate}

\textbf{2. Avance de un frente de luz en S'}
Un destello se produce en el origen común de los sistemas inerciales $S$ y $S'$ cuando ambos coinciden ($t= t'= 0$). A un tiempo $t$ la onda luminosa se ha expandido como una esfera. ¿Cuál es la coordenada que describe la ubicación del frente de onda del destello según $S'$?
\begin{enumerate}[a)]
    \item $x'= c/t'$
    \item $x'= \gamma x$
    \item $x'= ct+ vt$
    \item \textcolor{red}{$x'= ct'$}
\end{enumerate}

\textbf{3. Destello observado por S y S'}
Siguiendo con el mismo problema, consideremos la Transformación de Lorentz (TL) directa e inversa:
$$x'= \gamma(x - vt), \quad x= \gamma(x'+ vt')$$
Reemplazando las observaciones del pulso de luz hechas por $S$ y $S'$ llegamos a:
$$ct'= \gamma(ct - vt), \quad ct= \gamma(ct'+ vt')$$
Eliminando $t, t'$ podemos deducir:
\begin{enumerate}[a)]
    \item $\gamma= 1$
    \item $\gamma=(1 - \beta^2)^{-1}$
    \item $\gamma^2= v^2/c^2$
    \item \textcolor{red}{$\gamma^2= \frac{c^2}{c^2-v^2}$}
\end{enumerate}

\textbf{4. Propiedades de $\gamma$}
¿Cuál de estas propiedades obedece el factor de Lorentz $\gamma$?
\begin{enumerate}[a)]
    \item \textcolor{red}{$\gamma - \gamma^{-1}= \beta^2\gamma$}
    \item $\gamma< 1$
    \item $\gamma \approx \beta$ para $v \ll c$
    \item $\gamma \approx c$ para $v \approx c$
\end{enumerate}

\textbf{5. Transformación de Lorentz para el tiempo}
Consideremos de nuevo la TL directa e inversa escritas en 4-vectores:
$$x'_1= \gamma(x_1 - \beta x_0), \quad x_1= \gamma(x'_1+ \beta x'_0)$$
Sustituyendo la primera en la segunda, podemos obtener:
\begin{enumerate}[a)]
    \item $\beta x'_0= x_1 - \beta x_0$
    \item \textcolor{red}{$\beta x'_0= -\beta^2\gamma x_1+ \beta\gamma x_0$}
    \item $\beta x'_0= \gamma x_1+ \beta\gamma x_0$
    \item $\beta x'_0= \beta x_0$
\end{enumerate}

\textbf{6. Transformaciones de Lorentz}
En el sistema $S$ un destello ocurre en $x= 100$ km, $y= 10$ km, $z= 1$ km, $t= 0.5$ ms. El sistema $S'$ se mueve con velocidad $v= -0.8c$ a lo largo del eje común $xx'$. ¿Cuáles son las coordenadas del evento en $S'$?
\begin{enumerate}[a)]
    \item $x'= 100$ km, $y'= 10$ km, $z'= 1$ km, $t'= 0.5$ ms
    \item \textcolor{red}{$x' \approx 366.5$ km, $y'= 10$ km, $z'= 1$ km, $t' \approx 1.278$ ms}
    \item $x' \approx 83.3$ km, $y'= 10$ km, $z'= 1$ km, $t' \approx 0.3$ ms
    \item $x' \approx 120$ km, $y'= 12$ km, $z'= 2$ km, $t' \approx 0.8$ ms
\end{enumerate}

\textbf{7. Transformaciones de Lorentz}
En el sistema $S$ ocurre una explosión en $x= 5$ km, $t= 1$ ms. Para el observador $S'$ el evento está en $x'= 35.354$ km. Encuentre la velocidad relativa entre $S$ y $S'$ y el tiempo del evento medido en $S'$.
\begin{enumerate}[a)]
    \item $v=+0.1007 c$, $t'= 0.993$ ms
    \item $v= 0$, $t'= 1.000$ ms
    \item $v= -0.800 c$, $t'= 1.278$ ms
    \item \textcolor{red}{$v= -0.10065 c$, $t'= 1.0068$ ms}
\end{enumerate}
